% Source-derived chapter 10: Properties of the global stratification
% Original source: intersection-complexes-unramified-l-factors
\chapter{Properties of the global stratification}
\label{intersection-complexes-unramified-l-factors-chapter-10}
\source{intersection-complexes-unramified-l-factors}

\begin{remark}[Source section]
\label{intersection-complexes-unramified-l-factors-chapter-10-source}
\source{intersection-complexes-unramified-l-factors}
\source{intersection-complexes-unramified-l-factors}
This chapter reproduces the corresponding section of the retrieved
LaTeX source, including its definitions, statements, examples, and
proofs. It has not yet been translated into Lean.
\notready
\end{remark}

% The source section title was: Properties of the global stratification
\label{appendix:A}
\source{intersection-complexes-unramified-l-factors}

\subsection{The factorizable space of formal loops} \label{sect:multijet}
\source{intersection-complexes-unramified-l-factors}

We briskly review the definitions of multi-point versions of the 
spaces of formal arcs and formal loops. We refer the reader to \cite{KV}, \cite[\S 3.1]{Xinwen} for a more complete account.  

Let $C$ be a smooth curve over $k$. 
For any $N \in \mbb N$, we have the $N$th symmetric product
$C^{(N)}$ of $C$, which identifies with the Hilbert scheme $\on{Hilb}^N(C)$
parametrizing relative effective divisors in $C$ of degree $N$.

Recall that if $S$ is an affine scheme and $D \subset C\times S$ 
is a closed affine subscheme, we denote by $\wh C'_D$ the 
spectrum of the ring of regular functions
on the formal completion of $C\times S$ along $D$ (so $\wh C'_D$ is a true scheme, not merely a formal scheme). 
Let $\wh C^\circ_D := \wh C'_D \sm D$ denote
the open subscheme. 

For any $k$-scheme $X$, we define the global space of formal arcs
by the functor
\begin{align*} 
(\Scr L^+ X)_{C^{(N)}}(S) & = \{ D \in C^{(N)}(S), \gamma \in X(\wh C'_D) \}. 
\end{align*}
for affine test schemes $S$. 
By \cite[Proposition 2.4.1]{KV}, the functor $(\Scr L^+ X)_{C^{(N)}}$ is representable by a scheme of infinite type over $C^{(N)}$.
If $X$ is affine, then $(\Scr L^+ X)_{C^{(N)}}$ is affine
over $C^{(N)}$, cf.~\cite[2.4.3]{KV}. 
More specifically, 
define the space of $n$-jets $(\Scr L^+_n X)_{C^{(N)}}$ by 
\[ (\Scr L^+_n X)_{C^{(N)}}(S) = \{ D\in C^{(N)}(S), \gamma \in X(\wh C^n_D) \}, \]
where $\wh C^n_D$ denotes the $n$th infinitesimal neighborhood of $D$
in $C\times S$. Then $(\Scr L^+_n X)_{C^{(N)}}$ is representable
by a scheme over $C^{(N)}$, which is affine (resp.~of finite type) if $X$ is.  As $n$ varies 
the schemes $(\Scr L^+_n X)_{C^{(N)}}$ form a projective system of schemes 
with affine transition maps, and 
$(\Scr L^+_n X)_{C^{(N)}}$ is equal to the projective limit of this system.
If $X$ is smooth, the schemes $(\Scr L^+_n X)_{C^{(N)}}$ are
smooth over $C^{(N)}$ with smooth surjective transition maps (cf.~\cite[Lemma 2.5.1]{cpsii}).


We can also define the functor for the global loop space by 
\begin{align*} 
(\Scr L X)_{C^{(N)}}(S) & = \{ D \in C^{(N)}(S), \gamma \in X(\wh C^\circ_D) \}. 
\end{align*}
\emph{If $X$ is affine}, then $(\Scr L X)_{C^{(N)}}$ 
is representable by an ind-scheme ind-affine over $C^{(N)}$,
cf.~\cite[Proposition 2.5.2]{KV}. 
We have a closed embedding $(\Scr L^+ X)_{C^{(N)}} \into (\Scr L X)_{C^{(N)}}$.

\subsubsection{}\label{sect:Llocalization}
\source{intersection-complexes-unramified-l-factors}
Let $\Aut^0(k\tbrac t)$ denote the functor sending a $k$-algebra $R$
to the group of $R$-algebra automorphisms of $R\tbrac t$ that reduce to
the identity map mod $t$. Then $\Aut^0(k\tbrac t)$ is representable 
by the group scheme $\Spec k[a_1^{\pm 1}, a_2, a_3, \dotsc]$, cf.~\cite[(1.3.13)]{Xinwen}. 

\smallskip

There is an $\Aut^0(k\tbrac t)$-torsor $\on{Coord}^0(C) \to C$
classifying $v \in C$ together with an isomorphism $k\tbrac t \cong \mf o_v$
that sends $t$ to a uniformizer, cf.~\cite[(3.1.10)]{Xinwen}. 
We can think of $(\Scr L^+ X)_C, (\Scr L X)_C$ as twisted products 
\[
    (\Scr L^+ X)_C = C \ttimes \msf L^+ X, \qquad (\Scr L X)_C = C\ttimes \msf L X, 
\]
where $C \ttimes \msf L^+ X := \on{Coord}^0(C) \times^{\Aut^0(k\tbrac t)} \msf L^+ X$.


\begin{rem}
\source{intersection-complexes-unramified-l-factors}
\notready
The space $\Scr L X$ really lives over
the Ran space of $C$. Essentially this just means $\Scr L X$
only cares about the support of the divisor $D$ and not its multiplicities. 
More specifically for any finite set $I$ we have a map 
$C^I \to C^{(\abs I)}$ where $\abs I$ denotes the cardinality of $I$. 
Then the spaces $(\Scr L X)_{C^I} := C^I \times_{C^{(\abs I)}} (\Scr L X)_{C^{(\abs I)}}$ 
have a factorization monoid structure as
defined in \cite[Definition 2.2.1]{KV}, and we can think 
of the collection of these spaces as $(\Scr L X)_{\on{Ran}_C}$. 
This is certainly the more philosophically correct approach to 
considering the \emph{loop} space, but for technical simplicity 
it will suffice for our study of \emph{arc} spaces to work with $\Scr L^+ X$
over $\Sym C$.
\end{rem}

\subsubsection{} \label{sect:localGr}
\source{intersection-complexes-unramified-l-factors}
We can apply the constructions above to the algebraic group $G$. 
Since $G$ is smooth, $\Scr L^+ G$ is a group scheme formally smooth over 
$\Sym C$.

Consider the (factorizable) Beilinson--Drinfeld affine Grassmannian $\Gr_{G,\Sym C}$ defined in \S\ref{sect:YGr}. 
By Beauville--Laszlo's theorem (see \cite[Remark 2.3.7]{BD}, \cite[Proposition 3.1.9]{Xinwen}), we have an isomorphism 
$ \Gr_{G,\Sym C} \cong \Scr L G / \Scr L^+ G$. 


\subsubsection{} Let $X$ be an affine spherical 
$G$-variety. Define 
\[
    \Scr L^\bullet X := \Scr L X \sm \Scr L(X\sm X^\bullet),
\]
which admits an open embedding into $\Scr L X$. 
The $G$-action on $X$ induces a natural action of $\Scr L G$ on 
$\Scr L X$ and the subspace $\Scr L^\bullet X$ (resp.~$\Scr L^+ X$) is stable under the action of 
$\Scr L G$ (resp.~$\Scr L^+ G$). 

We will primarily be concerned with the global space of 
non-degenerate arcs
\[ (\Scr L^+X)^\bullet := \Scr L^+X \times_{\Scr L X} \Scr L^\bullet X.\] 
The study of the loop space $\Scr L^\bullet X$ is beyond the scope of 
this paper.


\subsection{Multi-point orbits} \label{sect:defstrata}
\source{intersection-complexes-unramified-l-factors}
We now consider the $\Scr L^+G$-orbits on $(\Scr L^+X)^\bullet$, and
we will prove a multi-point version of Proposition~\ref{prop:smoothfiberstrata}.

\subsubsection{What is going on at the level of $k$-points} 
A $k$-point of $\Ran_C$ is a nonempty finite subset 
$\{ v_i \in \abs C \}_{i\in I}$ of points on $C$. 
Then a $k$-point of $(\Scr L^+X)^\bullet_{\Ran_C}$ over $\{v_i\}$ consists
of points $x_i \in X(\mf o_{v_i}) \cap X^\bullet(F_{v_i})$. 
Each $x_i$ belongs to an orbit $X^\bullet(F_{v_i})_{G:\ch\theta_i}$
for a unique $\ch\theta_i \in \mathfrak c_X^-$ by Theorem~\ref{thm:Gorbits}(ii). 
The collection $\ch\theta_{i},\, i\in I$ forms a multiset
in $\mathfrak c_X^-$. 
Note that $\ch\theta_{i}$ may be zero. 
Therefore, the tuple $(x_i)_{i\in I}$ is a $k$-point in the product of orbits 
$\underset{i\in I}\prod \msf L^{\ch\theta_i}_{v_i} X$.

The idea for what follows is that this product of orbits only depends on
the unordered multiset $\{\ch\theta_i\}$, counted with multiplicity. 
Moreover since $\Cal C_0$ is a strictly convex cone, 
the set of formal sums $\sum_I \ch\theta_i \cdot v_i$ admits a positive grading.

\subsubsection{Construction}\label{sect:constructstrata}
\source{intersection-complexes-unramified-l-factors}
Let $\ch\Theta_0$ denote an (unordered) multiset in $\mathfrak c_X^-$, by 
which we mean a formal sum 
\[ \sum_{\ch\theta\in \mathfrak c_X^- } N_{\ch\theta} [\ch\theta] \in \Sym^\infty(\mathfrak c_X^-) \]
where all but finitely many $N_{\ch\theta}=0$. 
Note that we include the case $\ch\theta=0$ and $N_0$ may be nonzero. 
See \S\ref{def:partition} for notation.

We have an \'etale map 
\[ \oo C^{\abs{\ch\Theta_0}} = \underset{\ch\theta \in \mf c_X^-}{\oo\prod} \oo C^{N_{\ch\theta}} 
\to \underset{\ch\theta \in \mf c_X^-}{\oo\prod} \oo C^{(N_{\ch\theta})} = \oo C^{\ch\Theta_0}. 
\]
By the factorization property, the base change 
$\oo C^{\abs{\ch\Theta_0}} \xt_{\Sym C} \Scr L^+ X$ identifies with the $\abs{\ch\Theta_0}$-fold 
disjoint product of $C \ttimes \msf L^+ X$. 
By Proposition~\ref{prop:smoothfiberstrata}, we have the locally closed 
subscheme 
\begin{equation} \label{e:productstrata}
\source{intersection-complexes-unramified-l-factors}
\underset{\ch\theta \in \mf c_X^-}{\oo\prod} (C \ttimes \msf L^{\ch\theta} X)^{\oo\xt N_{\ch\theta}} 
\into \oo C^{\abs{\ch\Theta_0}} \xt_{\Sym C} \Scr L^+ X.
\end{equation}
This is stable under the action of $\prod_{\ch\theta} {\mf S}_{N_{\ch\theta}}$,
and therefore \eqref{e:productstrata} descends to 
a locally closed subscheme 
\[ \Scr L^{\ch\Theta_0} X \into 
\oo C^{\ch\Theta_0} \underset{\Sym C}\times \Scr L^+ X =: (\Scr L^+ X)_{\oo C^{\ch\Theta_0}}. \]

\begin{prop}
\source{intersection-complexes-unramified-l-factors}
\notready \label{prop:smoothstrata}
\source{intersection-complexes-unramified-l-factors}
The scheme $\Scr L^{\ch\Theta_0} X$ is formally smooth over $\oo C^{\ch\Theta_0}$,
and second projection induces a locally closed embedding $\Scr L^{\ch\Theta_0}X \into 
\Scr L^+ X$. 
\end{prop}
\begin{proof}
Formal smoothness follows from Proposition \ref{prop:smoothfiberstrata}. 
It remains to show that the second projection 
$\on{pr}_2 : \Scr L^{\ch\Theta_0} X \to \Scr L^+ X$
is a locally closed embedding. 
Note that $\on{pr}_2$ is the composition of the finite \'etale map 
\[ \on{pr}'_2: (\Scr L^+ X)_{\oo C^{\ch\Theta_0}} \to 
(\Scr L^+ X)_{\oo C^{( \abs{\ch\Theta_0})}} := \oo C^{(\abs{\ch\Theta_0})} \underset{\Sym C}\times \Scr L^+ X \]
and the open embedding $(\Scr L^+ X)_{\oo C^{(\abs{\ch\Theta_0})}} \into \Scr L^+ X$. 
Hence it suffices to show that the restriction 
$\Scr L^{\ch\Theta_0} X  \subset (\Scr L^+ X)_{\oo C^{\ch\Theta_0}} \to 
(\Scr L^+ X)_{\oo C^{(\abs{\ch\Theta_0})}}$
is locally closed.  
Let $Y$ be the closure of $\Scr L^{\ch\Theta_0} X$ in $(\Scr L^+ X)_{\oo C^{\ch\Theta_0}}$, 
so $Y'=\on{pr}'_2(Y)$ is 
a closed subscheme of $(\Scr L^+ X)_{\oo C^{(\abs{\ch\Theta_0})}}$
and $\on{pr}'_2(\Scr L^{\ch\Theta_0} X)$ is open in $Y'$. 
Observe that the induced map 
\begin{equation} \label{e:productstratapr2} 
\source{intersection-complexes-unramified-l-factors}
\Scr L^{\ch\Theta_0} X \to \on{pr}'_2(\Scr L^{\ch\Theta_0} X)
\end{equation}
is a bijection on geometric points: 
a point in the right hand side consists of an unordered set of points $\{ v_i \} \subset C$
and $x_i \in \msf L^+_{v_i} X$ such that if $\ch\theta_i$ denotes the $G$-valuation
of $x_i$, then $\sum_i [\ch\theta_i] = \ch\Theta_0$. 
There is a unique way to partition the $v_i$'s according to distinct values of $\ch\theta_i$'s
to get a point in $\Scr L^{\ch\Theta_0} X$. 
Therefore \eqref{e:productstratapr2} is \'etale and a bijection on geometric points, hence
an isomorphism. 
\end{proof}

\subsection{Proof of Lemma {\ref{lem:globstrata}}}
\label{proof:globstrata}
\source{intersection-complexes-unramified-l-factors}

Assume that $C$ is complete. 
Let $(\sM_X \xt \Sym C)^\bullet$ denote the substack of $\sM_X \xt\Sym C$
consisting of those pairs $(f,D)$ where $f(C \sm D) \subset X^\bullet/G$, i.e., 
$C\sm D$ is contained in the non-degenerate locus.
This is an open substack since $C$ is complete. 
Define the map
\begin{equation} \label{e:MtoLX}
\source{intersection-complexes-unramified-l-factors}
    (\Scr M_X \times \Sym C)^\bullet \to \Scr L^+ X / \Scr L^+G
\end{equation}
over $\Sym C$ by sending $(f,D)$ to $(f|_{\wh C'_D}, D)$. 
Here we are using the fact that an $\Scr L^+G$-torsor on an affine scheme $S$
is the same as a $G$-torsor on $\wh C'_D$ by formal lifting. 


Recall that we defined a partition $\ch\Theta$, or \emph{unordered multiset without zero}, in $\mathfrak c_X^-$, 
to mean an element of $\Sym^\infty(\mf c_X^- \sm 0)$. 
To such a partition $\ch\Theta$, we have a corresponding 
locally closed subscheme $\Scr L^{\ch\Theta} X \into \Scr L^+ X$
by Proposition~\ref{prop:smoothstrata}. 
Then we can \emph{define} $\Scr M^{\ch\Theta}_X$ 
to be the preimage of $\Scr L^{\ch\Theta} X / \Scr L^+G$ under the map 
\eqref{e:MtoLX}. 
One can check from the construction that 
this definition gives the description on $k$-points 
from \S\ref{sect:globstrat}.
By base change $\Scr M^{\ch\Theta}_X$ is a locally closed
subscheme of $\Scr M_X \times C^{(\abs{\ch\Theta})}$. 
In particular, $\Scr M^{\ch\Theta}_X$ is an algebraic 
stack locally of finite type over $k$. 

\begin{lem}
\source{intersection-complexes-unramified-l-factors}
\notready \label{lem:MtoLsm}
\source{intersection-complexes-unramified-l-factors}
The natural map $\Scr M^{\ch\Theta}_X \to \Scr L^{\ch\Theta} X / \Scr L^+G$ is formally 
smooth. 
\end{lem}
\begin{proof}
Let $S \into S'$ be a nilpotent thickening of affine schemes. 
Let $(f,D) \in \Scr M_X^{\ch\Theta}(S)$. 
This maps to the point $(f|_{\wh C'_D}, D) \in \Scr L^{\ch\Theta} X(S)$. 
Suppose that we have a lift of this point to 
$(\hat f, D') \in \Scr L^{\ch\Theta} X(S')$, where 
$D'\subset C\times S'$ is a relative effective Cartier divisor
and $\hat f: \wh C'_{D'} \to X/G$ is equivalent to 
the datum of a $G$-bundle $\hat{\Scr P}'_G$ on $\wh C'_{D'}$
and a section $\hat \sigma' : \wh C'_{D'} \to X \times^G \hat{\Scr P}'_G$.

We would like to lift $(f,D)$ to an $S'$-point of $\Scr M_X^{\ch\Theta}$
that maps to $(\hat f,D')$.
The map $f : C\times S \to X/G$
consists of the datum of a $G$-bundle $\Scr P_G$ on $C\times S$
and a section $\sigma : C\times S \to X \times^G \Scr P_G$
satisfying the condition that $\sigma(C\times S \sm D) \subset X^\bullet \times^G \Scr P_G = (H\bs G)\times^G \Scr P_G$.
The restriction $\sigma|_{C \times S \sm D}$ 
gives a reduction of $\Scr P_G|_{C\times S\sm D} \cong G\times^H \Scr P_H$ to an
$H$-bundle $\Scr P_H$ on $C\times S \sm D$ such that
$\sigma|_{C\times S\sm D}$ identifies with
the canonical section 
\[ C\times S \sm D \cong H\bs \Scr P_H \into (H\bs G)\overset H\times \Scr P_H \cong (H\bs G) \overset G\times \Scr P_G|_{C\times S \sm D}\]
corresponding to $H1\in H\bs G$.

The obstruction to lifting $\Scr P_H$ to an $H$-bundle $\Scr P'_H$ on $C\times S' \sm D'$ is an element in $H^2(C\times S\sm D, \mf h_{\Scr P_H} \otimes_{\Scr O_S} I)$ where $I$
is the zero ideal of $S\into S'$ and $\mf h_{\Scr P_H}$ denotes the quasicoherent sheaf on $C\times S \sm D$ obtained
by twisting the adjoint representation of $H$ by $\Scr P_H$.
This obstruction vanishes since $C\times S \sm D$ has
relative dimension $1$ over $S$ and we can compute
cohomology over the Zariski site. 
Thus, we obtained an $H$-bundle $\Scr P'_H$ over $C\times S' \sm D'$.
Let $\sigma' : C\times S'\sm D' \to X \times^G (G\times^H\Scr P'_H)$ denote the corresponding section.

We know that after base change to $S$, there exists an isomorphism
\[ \tau : \hat{\Scr P}'_G|_{\wh C^\circ_D} \cong \Scr P_G|_{\wh C^\circ_D} \cong (G\overset H\times \Scr P'_H)|_{\wh C^\circ_D} \]
such that $\tau \circ \hat\sigma'|_{\wh C^\circ_D} = \sigma'|_{\wh C^\circ_D}$. This is equivalent to a section 
$\beta: \wh C^\circ_D \to 
\Scr P'_H\times^H G\times^G \hat{\Scr P}'_G$
such that $\beta$ is sent under 
\begin{equation} \label{e:twistedGtoX} 
\source{intersection-complexes-unramified-l-factors}
    \Scr P'_H \overset H\times G \overset G\times \hat{\Scr P}'_G \to (C\times S')\times X \overset G\times \hat{\Scr P}'_G
\end{equation}
 to the restriction $\hat \sigma'|_{\wh C^\circ_{D}}$.
The map \eqref{e:twistedGtoX} is smooth since $G\to X^\bullet=H\bs G$ is smooth.
The scheme $\wh C^\circ_{D'}$ is affine since $S'$ is affine.
The zero ideal of $\wh C^\circ_D\into \wh C^\circ_{D'}$ is still
nilpotent, so by formal smoothness of \eqref{e:twistedGtoX}, we
can lift $\beta$ to a section $\beta'$ that maps to $\hat \sigma'|_{\wh C^\circ_{D'}}$.
Such a section $\beta'$ is equivalent to an isomorphism 
$\tau' : \hat{\Scr P}'_G|_{\wh C^\circ_{D'}} \cong (G\times^H \Scr P'_H)|_{\wh C^\circ_{D'}}$ such that $\tau' \circ \hat\sigma' |_{\wh C^\circ_{D'}} = \sigma'|_{\wh C^\circ_{D'}}$.
By Beauville--Laszlo's theorem (Lemma~\ref{lem:B-L}),
the data $((\hat{\Scr P}'_G,\hat\sigma'), (\Scr P'_H,\sigma'), \tau')$ descends to a map $f' : C\times S' \to X/G$. 
By construction, $(f',D')$ is an $S'$-point of $\Scr M^{\ch\Theta}_X$
lifting $f$.
\end{proof}


\begin{proof}[{Proof of Lemma~\ref{lem:globstrata}}] 
Lemma~\ref{lem:MtoLsm} and Proposition~\ref{prop:smoothstrata} together
imply that $\Scr M^{\ch\Theta}_X$ is formally smooth over $k$. 
Since $\Scr M^{\ch\Theta}_X$ is locally of finite type,
it is therefore smooth over $k$.

We claim that the first projection $\on{pr}_1 : \Scr M_X\times C^{(\abs{\ch\Theta})} \to \Scr M_X$ induces a locally closed
embedding $\Scr M^{\ch\Theta}_X \into \Scr M_X$. 
Let $Z \subset \Scr M_X \times C^{(\abs{\ch\Theta})}$ denote the substack
with $S$-points consisting of those $(f,D)$ such that 
$f^{-1}(X^\bullet/G)\cap D = \emptyset$. 
Since $D$ is faithfully flat over $S$, the image of $f^{-1}(X^\bullet/G)\cap D$ in $S$ is open. Therefore, $Z$
is a closed substack of $\Scr M_X \times C^{(\abs{\ch\Theta})}$.
Since $\ch\Theta$ is a multiset without zero, 
observe that $\Scr M^{\ch\Theta}_X$ embeds into 
$Z$. Now $(f,D)\in Z(k)$ satisfies the property that the support 
of $D\in C^{(\abs{\ch\Theta})}$ is contained in the support of $C \sm f^{-1}(X^\bullet/G)$. We deduce that ${\on{pr}_1}|_Z : Z \to \Scr M_X$ is 
proper and quasifinite, hence finite. 
On the other hand, $\on{pr}_1|_{\sM^{\ch\Theta}_X}$ is injective
on $k$-points, and 
$({\on{pr}_1}|_Z)^{-1}(\on{pr}_1(\Scr M^{\ch\Theta}_X)) = \Scr M^{\ch\Theta}_X$. 
From this we deduce that $\on{pr}_1|_{\sM^{\ch\Theta}_X}$ 
is a locally closed embedding.
\end{proof}



\subsection{Generic-Hecke modifications} 
We review the notion of \emph{generic-Hecke modifications} between quasimaps
introduced in \cite[\S 2.2]{GN}, applied to our situation. 
We assume that $B$ acts simply transitively on $X^\circ$ and that $H$ is connected. 

\subsubsection{Function-theoretic analog} 
We explain the idea behind the generic-Hecke modifications at the level of sets; 
this construction appears in the geometric Langlands
program in the construction of Whittaker models, cf.~\cite[\S 5.3.1]{G:outline}.

We use the notation of \S\ref{sect:adelic}.
For any finite subset $\underline v \subset \abs C$, 
let $\bbA^{\ul v} = \prod'_{v'\notin \ul v} F_{v'},\, 
\bbA_{\ul v} = \prod_{v'\in \ul v} F_{v'}$ and similarly for 
$\mbb O^{\ul v}, \mbb O_{\ul v}$. Then we can consider 
the set  
\begin{equation} \label{e:genericHeckeset} 
\source{intersection-complexes-unramified-l-factors}
 G(\bbA^{\ul v})/ G(\mbb O^{\ul v}) \xt H(\bbA_{\ul v})/ H(\mbb O_{\ul v})   
\end{equation}
which maps to $H(\Bbbk) \bs G(\bbA) / G(\mbb O)$ in two ways: 
(i) by projecting along the first factor to $( H(\Bbbk) \xt_{H(\bbA_{\ul v})} H(\mbb O_{\ul v})) \bs G(\bbA^{\ul v}) / G(\mbb O^{\ul v})
\subset H(\Bbbk) \bs G(\bbA) / G(\mbb O)$
and (ii) by projecting to 
\[ H(\Bbbk) \bs \Bigl( G(\bbA^{\ul v}) / G(\mbb O^{\ul v}) \xt H(\bbA_{\ul v})/H(\mbb O_{\ul v}) \Bigr) 
\subset H(\Bbbk) \bs G(\bbA)/ G(\mbb O).
\]
Every meromorphic quasimap (element of $H(\Bbbk) \bs G(\mbb A) / G(\mbb O)$) 
belongs to the image of the second projection for some 
$\underline v$. If $H$ is connected, then by \emph{weak approximation} 
$H(\bbA_{\underline v}) = H(\Bbbk) H(\mbb O_{\underline v})$ so 
every quasimap also belongs to the image of the first projection.
Thus, the union of \eqref{e:genericHeckeset} over all finite subsets $\ul v$ 
defines a groupoid acting on the set of quasimaps.


\subsubsection{}
We define the ind-stack $\Scr H_{H,\sM_X}$ of \emph{generic-Hecke modifications} 
to be the stack classifying data 
\[ (\Scr P^1_G, \Scr P^2_G, \sigma_1,\sigma_2; \ul v, \tau) \]
where $(\Scr P^i_G, \sigma_i) \in \sM_X$, 
$\ul v \in \Sym C$ is a divisor with support $\ul v$ contained in the non-degenerate locus
$\sigma_i^{-1}(X^\bullet \xt^G \sP^i_G)$ for both $i=1,2$
and $\tau$ is an isomorphism of $G$-bundles 
\[ \tau: \Scr P^1_G|_{C\sm {\ul v}} \cong \Scr P^2_G|_{C\sm {\ul v}} \]
such that the following diagram commutes
\[ 
\begin{tikzcd} 
C \sm {\ul v} \ar[r, "\sigma_1"] \ar[rd, swap,"\sigma_2"] & \Scr P^1_G \overset G\times X |_{C\sm {\ul v}} \ar[d, "\tau" ] \\
& \Scr P^2_G \overset G \times X |_{C\sm {\ul v}}. 
\end{tikzcd} 
\]
Note that the definition only depends on the support of $\ul v$ and not its multiplicities. 

We call a generic-Hecke modification \emph{trivial} if the isomorphism $\tau$
extends to an isomorphism over all of $C$. 
We have the natural projections 
\[ \Scr M_{X} \overset{h^\leftarrow}\leftarrow \Scr H_{H,\sM_X} \overset{h^\to}\to \Scr M_{X}, \] 
and $\Scr H_{H,\sM_X} \to \Sym C$. 
By definition, the generic-Hecke correspondence preserves the strata 
$\sM_X^{\ch\Theta}$. 

Define a \emph{smooth generic-Hecke correspondence} to be any stack $U$ equipped
with smooth maps 
\[ \Scr M_{X} \overset{h^\leftarrow_U}\leftarrow U \overset{h^\to_U}\to \Scr M_{X} \] 
such that there exists a map $U\to \Scr H_{H,\sM_X}$ such that
the following diagram commutes
\[ 
\begin{tikzcd} 
& U \ar[ld, swap, "h^\leftarrow_U" near start] \ar[d] \ar[rd, "h^\to_U" near start] \\
\Scr M_{X} & \Scr H_{H,\sM_X} \ar[l,swap,"h^\leftarrow" near start] \ar[r,"h^\to" near start] &
\Scr M_{X} 
\end{tikzcd}
\]
Call a smooth generic-Hecke correspondence $U$ trivial if the image of 
$U \to \Scr H_{H,\sM_X}$ consists of trivial generic-Hecke modifications. 

Define a morphism of smooth generic-Hecke correspondences to be a map 
$p : U_1 \to U_2$
such that $h^\leftarrow_{U_2}\circ p = h^\leftarrow_{U_1}$ and $h^\to_{U_2}\circ p = h^\to_{U_1}$.

\subsubsection{Generic-Hecke equivariant sheaves}
We define a \emph{generic-Hecke equivariant} perverse sheaf on $\sM_X$ 
to be an object $\sF \in \rmP(\sM_X)$ of the category of perverse sheaves on $\sM_X$
equipped with isomorphisms 
\[ \phi_U : h^{\leftarrow *}_U(\sF) \overset\sim\to h_U^{\rightarrow *}(\sF) 
\] 
for every smooth generic-Hecke correspondence $U$, satisfying some natural 
conditions (see \cite[\S 2.3]{GN} for details).

For any $\ch\Theta \in \Sym^\infty(\mf c_X^- \sm 0)$, 
the uniqueness of the IC complex endows $\IC_{\barM^{\ch\Theta}_X}$ 
with the structure of a generic-Hecke equivariant perverse sheaf. 
(In general, when $H$ is connected, the condition of generic-Hecke equivariance
is a \emph{property}, not additional structure, of a perverse sheaf on $\sM_X$
by \cite[Proposition 3.5.2]{GN}.)

\subsubsection{}
Fix $\ch\theta \in \mathfrak c_X^- \sm 0$. 
Recall that by definition $\sM_X^{\ch\theta}$ is a substack of $\sM_X \xt C$. 
For a fixed $v_0 \in \abs C$, define $\sM_{X,v_0}^{\ch\theta} := \sM_X^{\ch\theta} \xt_C v_0$
to be the \emph{based} stratum consisting of maps $C \to X/G$ such that $C\sm v_0$
maps to $H\bs \pt$ and the $G$-valuation at $v_0$ is $\ch\theta$.
Observe that the generic-Hecke modifications preserve the substack
$\sM_{X,v_0}^{\ch\theta}$. 

\begin{prop}[{\cite[Proposition 3.5.1]{GN}}]
\source{intersection-complexes-unramified-l-factors}
\notready 
\label{prop:genHecke-transitive}
\source{intersection-complexes-unramified-l-factors}
If $H$ is connected, then all geometric points of $\sM_{X,v_0}^{\ch\theta} \subset \sM_X$
are equivalent under the equivalence relation generated by 
the generic-Hecke correspondences. 
\end{prop}

The preimage of $C\subset \Sym C$ in $\Scr H_{H,\sM_X}$ may be realized as the
twisted product of an open substack of $\sM_X \xt C$ with the affine Grassmannian
$\Gr_H$. From this it is not hard to deduce that all geometric points of 
$\sM_{X,v_0}^{\ch\theta}$ are equivalent under \emph{smooth} generic-Hecke correspondences. 

\subsubsection{} 
Recall the $G$-Hecke action defined in Proposition~\ref{prop-const:action}. 
By construction, this action commutes with the generic $H$-Hecke modifications (in 
our set-theoretic description above, we are considering the action of right
$G(F_{v'})$-translation at some point $v' \ne v_0$):

\begin{lem}
\source{intersection-complexes-unramified-l-factors}
\notready \label{lem:act-equivariant}
\source{intersection-complexes-unramified-l-factors}
The map $\act_C : \Scr M_X \ttimes \ol\Gr^{\ch\theta}_{G,C} \to \Scr M_X \xt C$ is equivariant with respect to the generic-Hecke modifications
away from the marked point in $C$. 
\end{lem}




\subsubsection{Proof of Proposition \ref{prop:globwhitney}} \label{proof:whitney}
\source{intersection-complexes-unramified-l-factors}

Proposition~\ref{prop:genHecke-transitive} allows us to consider the based strata with 
generic-Hecke modifications
as an analog of the stratification by orbits of a group action. 
This will be the idea behind the proof of Proposition~\ref{prop:globwhitney}, together
with the following fact: 

\begin{thm}[{\cite[Theorem 2]{Kaloshin}}]
\source{intersection-complexes-unramified-l-factors}
\notready  \label{thm:K}
\source{intersection-complexes-unramified-l-factors}
Let $S$ be a smooth stratum in an algebraically stratified scheme of finite type over $k=\mbb C$. 
Then the subset of points in $S$ that do not satisfy Whitney's condition B 
form a constructible subset of dimension strictly lower than $\dim S$.
\end{thm}

Thus, if we show that any two points in a given stratum have neighborhoods in $\sM_X$ 
that are smooth locally isomorphic and compatible with the stratification, 
Theorem~\ref{thm:K} will imply that every point
must satisfy Whitney's condition B.

\smallskip

Fix a connected component of $\sM_X^{\ch\Theta}$ 
for some $\ch\Theta \in \Sym^\infty(\mathfrak c_X \sm 0)$. 
Since Whitney's condition B is local in the smooth topology, 
it suffices by Corollary~\ref{cor:Mcover} to show that every point in 
$\sY^{\ch\lambda,\ch\Theta}$
satisfies Whitney's condition for all $\ch\lambda \in \mathfrak c_X$.
Now the graded factorization property of $\sY$ allows us to 
to reduce to the case where $\ch\Theta = [\ch\theta]$ is singleton. 
By Proposition~\ref{prop:changecurve}, we may replace our curve $C$ 
with $\bbA^1 = \mbb P^1 \sm \infty$, i.e., $\sY^{\ch\lambda}=\sY^{\ch\lambda}(\bbA^1)$.
Now $\bbA^1$ acts on itself by translation, which induces
an $\bbA^1$-action by automorphisms on $\sY^{\ch\lambda}(\bbA^1)$. 
This action allows us to move the degenerate point of any 
$y \in \sY^{\ch\lambda,\ch\theta}(\bbA^1)$ to $v_0=0 \in \bbA^1$, i.e., 
the map $y:\bbA^1 \to X/B$ sends $\bbA^1 \sm 0 \to X^\bullet/B$ and 
the $G$-valuation at $v_0$ equals $\ch\theta$.
Thus, we are reduced to showing that any point in 
$\sY^{\ch\lambda,\ch\theta}$ with degenerate point at $v_0$ satisfies Whitney's condition.

Embed $\bbA^1 = \mbb P^1 \sm \infty \subset \mbb P^1$ so we 
also have an open embedding $\sY^{\ch\lambda}(\bbA^1) \subset \sY^{\ch\lambda}(\mbb P^1)$. 
Lemma~\ref{lem:localglobalyoga} shows that $\sY^{\ch\lambda}(\bbP^1)$ 
is smooth locally isomorphic to $\sM_X=\sM_X(\bbP^1)$ in a way that preserves strata 
and degenerate points. Therefore, we may reduce to checking
that any point in $\sM_{X,v_0}^{\ch\theta}$ satisfies Whitney's condition B.
Proposition~\ref{prop:genHecke-transitive} implies that
all such points are equivalent under smooth generic-Hecke correspondences 
(which preserve strata), so either they all satisfy or fail to satisfy
Whitney's condition. By Theorem~\ref{thm:K}, we deduce that every
point satisfies Whitney's condition B.

\medskip

The same argument as above also shows that the closure of any stratum in $\sM_X$
must equal a union of strata. 

\subsubsection{Arbitrary characteristic} 

Let the characteristic of $k$ be arbitrary. 

\begin{prop}
\source{intersection-complexes-unramified-l-factors}
\notready \label{prop:Whitney-poschar}
\source{intersection-complexes-unramified-l-factors}
For any $\ch\Theta \in \Sym^\infty(\mf c_X^-\sm 0)$,
the complex $\IC_{\barM^{\ch\Theta}_X}$ 
is constructible with respect to the fine stratification of $\sM_X$.
\end{prop}

\begin{proof}
The argument is exactly the same as in the proof of Proposition~\ref{prop:globwhitney} 
above, except that we replace the use of Theorem~\ref{thm:K} with
the fact that for any connected smooth stratum $S$ in $\sM_X$, 
there exists some open $U\subset S$ such that 
$\rmH^i(\IC_{\barM^{\ch\Theta}_X})|_U$ is a local system for all $i$.
We implicitly use the uniqueness of the IC complex, which in
particular implies that $\IC_{\barM^{\ch\Theta}_X}$ is generic-Hecke equivariant.
\end{proof}


\subsection{Proof of Theorem~\ref{thm:heckeIC}} \label{sect:pf:heckeIC}
\source{intersection-complexes-unramified-l-factors}

For $\ch\Theta \in \Sym^\infty(\mathfrak c_X^-\sm 0)$, 
let $\rmP_{\Scr L^+ G}(\ol\Gr^{\ch\Theta}_{G,C^{\ch\Theta}})$ denote the
category of perverse sheaves on $\ol\Gr^{\ch\Theta}_{G,C^{\ch\Theta}}$ that
are equivariant with respect to the action of $(\Scr L^+ G)_{C^{(\abs{\ch\Theta})}}$, 
considered as a group scheme over $\Sym C$ 
(defined in \S\ref{sect:multijet}). 
For $\Scr F \in \rmP(\Scr M_X)$ and 
$\Scr G \in \rmP_{\Scr L^+G}(\ol\Gr^{\ch\Theta}_{G,C^{\ch\Theta}})$, we 
can form the twisted external product 
\[ \Scr F \mathbin{\tilde\bt} \Scr G \in \rmP(\Scr M_X \ttimes \ol\Gr^{\ch\Theta}_{G,C^{\ch\Theta}}) \]
with respect to the projections of 
$\widehat{\Scr M_X} \xt^{\Scr L^+G} \ol\Gr^{\ch\Theta}_{G,C^{\ch\Theta}}$ to 
$\Scr M_X$ and $\Scr L^+G \bs \ol\Gr^{\ch\Theta}_{G,C^{\ch\Theta}}$.
Then we define the external convolution product by 
\[ \Scr F \star \Scr G := \act_{\sM,!}( \Scr F \mathbin{\tilde\bt} \Scr G ) \in \rmD^b_c(\Scr M_X). \]

We have introduced all the ingredients in the statement of Theorem \ref{thm:heckeIC}. The remainder of this section will be devoted to its proof. 




Observe that $\IC_{\barM_X^0} \mathbin{\tilde\bt} 
\IC_{\ol\Gr^{\ch\Theta}_{G,C}}
= \IC_{\barM^0_X \ttimes \ol\Gr^{\ch\Theta}_{G,C}}$, which we 
will simply denote $\Scr{IC}$ for brevity. 
We have a stratification 
\[ \barM_X^0 \ttimes\ol\Gr^{\ch\Theta}_{G,C^{\ch\Theta}} 
= \bigcup_{\ch\Theta',\ch\Theta''} 
 \barM_X^{\ch\Theta'} \ttimes \Gr^{\ch\Theta''}_{G,\oo C^{\ch\Theta''}} \] 
running over all $\ch\Theta' \succeq 0$ 
and $\ch\Theta'' \ge \ch\Theta_1$ where $\ch\Theta$ refines $\ch\Theta_1$ (the definition of $\ch\Theta'' \ge \ch\Theta_1$ is analogous to that of $\succeq$).  
Lemma~\ref{lem:actimage}(i) implies that $\act_\sM^{-1}(\sM^{\ch\Theta}_X)$
is contained in the dense open stratum corresponding to 
$\ch\Theta'=0$ and $\ch\Theta'' = \ch\Theta$.


\subsubsection{} \label{sect:ICglobineq}
\source{intersection-complexes-unramified-l-factors}
By Proposition~\ref{prop:globwhitney} we know that 
$\IC_{\barM_X^0}$ is constructible with respect to the fine stratification of 
$\sM_X$. Therefore, $\Scr{IC}$ is constructible
with respect to the stratification above, so 
if we let $L_{\mathrm{glob}}^{\ch\Theta', \ch\Theta''}$ denote
the $*$-restriction of $\Scr{IC}$ to the stratum 
$\sM_X^{\ch\Theta'} \ttimes \Gr^{\ch\Theta''}_{G,\oo C^{\ch\Theta''}}$, we have that its cohomology sheaves are local systems. 
Since $L^{\ch\Theta', \ch\Theta''}_{\mathrm{glob}}$ is the 
restriction of an IC complex, it lives in perverse cohomological degrees 
$\le 0$, and the inequality is strict unless $\ch\Theta'=0$ and $\ch\Theta''=\ch\Theta$. 
Since its cohomology sheaves are local systems, this implies that 
$L^{\ch\Theta', \ch\Theta''}_{\mathrm{glob}}$ 
lives in usual cohomological degrees 
$\le -\dim( \sM_X^{\ch\Theta'} \ttimes \Gr^{\ch\Theta''}_{G,\oo C^{\ch\Theta''}})$, and the inequality is strict unless 
$\ch\Theta'=0$ and $\ch\Theta''=\ch\Theta$. 
We are abusing notation here: $\sM^{\ch\Theta'}_X$ may not be connected, in which case each connected component should be considered separately. 
 
\smallskip

We abuse notation and also use $L^{\ch\Theta',\ch\Theta''}_{\mathrm{glob}}$
to denote its $!$-extension to $\sM_X \ttimes \ol\Gr^{\ch\Theta}_{G, C^{\ch\Theta}}$. 
Then by the characterizing properties of the intermediate extension 
(and the fact that $\Scr{IC}$ is Verdier self-dual), 
Theorem~\ref{thm:heckeIC} is equivalent
to the following assertion: 

\begin{prop}
\source{intersection-complexes-unramified-l-factors}
\notready \label{prop:heckeIC}
\source{intersection-complexes-unramified-l-factors}
For $\ch\Theta',\ch\Theta''$ as above, 
consider the complex $\act_{\sM,!}(L^{\ch\Theta',\ch\Theta''}_{\mathrm{glob}})$.  
Then:
\begin{enumerate}
\item It lives in $^p\rmD^{\le -1}(\sM_X)$ unless $\ch\Theta'=0$ and $\ch\Theta''=\ch\Theta$.
\item The $*$-restriction of 
$\act_{\sM,!}(L^{0,\ch\Theta}_{\mathrm{glob}})$ to $\barM^{\ch\Theta}_X\sm \sM^{\ch\Theta}_X$ lives in perverse cohomological degrees $\le -1$.
\item There is a natural identification
$\act_{\sM,!}(L^{0,\ch\Theta}_{\mathrm{glob}})|^*_{\sM^{\ch\Theta}_X} \cong \IC_{\sM^{\ch\Theta}_X}$. 
\end{enumerate}
\end{prop}

Point (iii) follows immediately from Theorem \ref{thm:comp}(iii). 
Points (i)-(ii) concern the cohomological degrees 
of the $*$-restriction of $\act_{\sM,!}(L^{\ch\Theta',\ch\Theta''}_{\mathrm{glob}})$
to a stratum of $\barM^{\ch\Theta}_X \sm \sM_X^{\ch\Theta}$. 
To simplify notation we only consider the restriction
to a stratum of the form $\sM_X^{\ch\eta}$ where $\ch\eta \in \mf c_X^- \sm 0$. 
The general case is proved in the same way by considering multiple points on $C$ 
simultaneously. 

\smallskip

By Theorem \ref{thm:comp} and Proposition \ref{prop:closure-rel}, the preimage 
$\act_\sM^{-1}(\sM_X^{\ch\eta})$ intersects $\sM_X^{\ch\Theta'} \ttimes \Gr^{\ch\Theta''}_{G,\oo C^{\ch\Theta''}}$ only if 
$\ch\Theta' = [\ch\theta']$ and $\ch\Theta'' = [\ch\theta'']$ are singletons 
($\ch\theta'$ is allowed to be $0$). 
Moreover we must have $\ch\theta'' \ge \deg(\ch\Theta)$ and $\ch\eta \succeq \ch\theta' + \ch\theta''$. 

\subsubsection{Fiber dimension}
Fix a point $v \in \abs C$ and consider the subscheme $\sM_{X,v}$
of maps where there is only one $G$-degenerate point at $v$. 
Let $\sM_{X,v}^{\ch\eta}$ denote the substack of $\sM_{X,v}$ where the $G$-valuation at
this degenerate point is $\ch\eta$.
Then the preimage $\act_\sM^{-1}(\sM_{X,v}^{\ch\eta}) \cap (\sM_X^{\ch\theta'} \ttimes \Gr^{\ch\theta''}_{G,C})$ is contained in $\sM_X^{\ch\theta'} \ttimes \Gr^{\ch\theta''}_{G,v}$.
The restricted map 
\begin{equation} \label{e:actrestricted}
\source{intersection-complexes-unramified-l-factors}
 \act_\sM : \sM_X^{\ch\theta'} \ttimes \Gr^{\ch\theta''}_{G,v} \to \sM_{X,v} 
\end{equation}
is equivariant with respect to generic-Hecke modifications away from $v$ by Lemma \ref{lem:act-equivariant}. Meanwhile these modifications act transitively on the stratum 
$\sM_{X,v}^{\ch\eta}$. 
Therefore we deduce that the fibers of \eqref{e:actrestricted} over 
all points in $\sM_{X,v}^{\ch\eta}$  are isomorphic to one another. 
Let this fiber dimension be denoted $d$. 

We have a special point $\{ t^{\ch\eta} \} = \msf Y^{\ch\eta,\ch\eta}_v \to \sM^{\ch\eta}_{X,v}$
by Corollary~\ref{cor:Ytheta}. 
We deduce from Proposition~\ref{prop:Yact-strat} that 
the preimage of $\msf Y^{\ch\eta}_v$ under the map \eqref{e:actrestricted} has a 
stratification by 
\[ \msf Y^{\ch\eta - \ch\nu, \ch\theta'} \ttimes (\msf S^{\ch\nu} \cap \Gr^{\ch\theta''}_G) \]
where $\ch\nu$ runs over the weights of $V^{\ch\theta''}$.
Therefore
\begin{equation} \label{e:actfiberdim}
\source{intersection-complexes-unramified-l-factors}
d \le \underset{\ch\nu}{\on{max}} (\dim \msf Y^{\ch\eta-\ch\nu,\ch\theta'} + \brac{\rho_G,\ch\nu-\ch\theta''}). 
\end{equation}


\subsubsection{Passage to Zastava model}
By Corollary~\ref{cor:Mcover}, it suffices to prove 
Proposition~\ref{prop:heckeIC} after base change to the Zastava model.
Consider the fiber product diagram 
\[\begin{tikzcd}
Z^{?,\ch\lambda,\ch\theta',\ch\theta''} \ar[r] \ar[d] & \sY^{\ch\lambda} \ar[d] \\
\sM_X^{\ch\theta'} \ttimes \Gr^{\ch\theta''}_{G,C} \ar[r, "\act_\sM"] & 
\sM_X 
\end{tikzcd}
\]
which is the analog of \eqref{e:Z0act-diagram} where we allow $v\in \abs C$ to vary. 
We pull back along this diagram for all $\ch\lambda$ large enough. 
The corresponding diagram for strata is given by \eqref{e:Z0act-diagram}.
We $*$-pullback with a shift by the fiber dimension of 
$\sY^{\ch\lambda}\to \sM_X$. 
With this shift, the discussion of \S\ref{sect:ICglobineq} 
implies that $L^{\ch\Theta',\ch\eta}_\glob$ goes to 
a complex 
\[ L^{?,\ch\lambda,\ch\theta',\ch\theta''}_\flag \in \rmD^b_c(Z^{?,\ch\lambda,\ch\theta',\ch\theta''}), \]
which lives in usual cohomological degrees 
$\le -\dim( Z^{?,\ch\lambda,\ch\theta',\ch\theta''} )$, and
the inequality is strict unless $\ch\theta'=0$ and $\ch\theta''=\ch\theta$. 
The usual cohomology sheaves of $L^{?,\ch\lambda,\ch\theta',\ch\theta''}_{\mathrm{flag}}$ 
are local systems. 

To prove Proposition~\ref{prop:heckeIC} it is
enough, by the definition of the perverse t-structure,
to prove the following:

\begin{lem}
\source{intersection-complexes-unramified-l-factors}
\notready \label{lem:heckeICflag}
\source{intersection-complexes-unramified-l-factors}
Let $\ch\theta',\ch\theta'',\ch\eta,\ch\lambda,d$ be as above. 
Then we have 
\[
    -\dim( Z^{?,\ch\lambda,\ch\theta',\ch\theta''} ) + 2 d \le - \dim \sY^{\ch\lambda,\ch\eta} 
\]
and the inequality is strict if $\ch\theta' = 0$ and $\ch\theta'' \ne \ch\eta$. 
\end{lem}
Note that the statement of the lemma has nothing to do with $\ch\theta$.

\begin{proof}
By \S\ref{sect:freemonoid}, we may assume that $\mf c_{X^\bullet} = \mbb N^{\Cal D}$. 
Now it makes sense to talk about $\len(\ch\lambda)$ for $\ch\lambda \succeq 0$, cf.~\S\ref{sect:Y0cc}. 
By Lemma \ref{lem:oYconnected} and Corollary~\ref{cor:Ycomp2},
we have that $\dim \sY^{\ch\lambda,\ch\eta} = \len(\ch\lambda-\ch\eta)+1$. 

For $\ch\nu$ a weight of $V^{\ch\theta''}$, we have 
$\dim \msf Y^{\ch\eta-\ch\nu,\ch\theta'} \le \frac 1 2(\dim \sY^{\ch\eta-\ch\nu,\ch\theta'}-1)$
by Proposition~\ref{prop:centralstrata}, unless $\ch\eta=\ch\nu$.  
If $\ch\theta'=0$, then $\dim \sY^{\ch\eta-\ch\nu,0} = \len(\ch\eta-\ch\nu)$. 
Otherwise $\dim \sY^{\ch\eta-\ch\nu,\ch\theta'} = \len(\ch\eta-\ch\nu)+1$. 
We also have $\brac{\rho_G,\ch\nu-\ch\theta''} = \frac 1 2 \len(\ch\nu-\ch\theta'')$. 
Therefore \eqref{e:actfiberdim} implies that 
\begin{equation} \label{ineq:dfiber}
\source{intersection-complexes-unramified-l-factors}
 d \le \frac 1 2 \len(\ch\eta-\ch\theta'-\ch\theta'')
\end{equation}
and the inequality is strict if $\ch\theta' = 0$ and $\ch\eta$ is not a weight of $V^{\ch\theta''}$. 

In the case $\ch\theta'=0$ and $\ch\eta$ is a weight of $V^{\ch\theta''}$ not equal to 
$\ch\theta''$, we claim the inequality \eqref{ineq:dfiber} above is still strict. 
Indeed, equality can only hold if $\ch\eta=\ch\nu$ and an open subvariety of 
$\msf Y^{0,0} \ttimes (\msf S^{\ch\eta}\cap \Gr^{\ch\theta''}_G)$ is sent to 
the special point $t^{\ch\eta}$ under \eqref{e:actrestricted}. 
However we know from Lemma~\ref{lem:MVcentralfiber} that every irreducible
component of $\msf S^{\ch\eta}\cap \Gr^{\ch\theta''}_G$ generically maps to 
the stratum $\sM^{\ch\theta''}_G$. Thus if $\ch\eta \ne \ch\theta''$, 
the inequality \eqref{ineq:dfiber} is strict.


Thus to prove the lemma it suffices to show that 
\[
 \len(\ch\lambda-\ch\theta'-\ch\theta'')+1 = \len(\ch\eta-\ch\theta'-\ch\theta'') + \len(\ch\lambda-\ch\eta) + 1 \le 
\dim (Z^{?,\ch\lambda,\ch\theta',\ch\theta''}).
\]
Proposition~\ref{prop:Yact-strat} (twisted by $C$) implies that 
$Z^{?,\ch\lambda,\ch\theta',\ch\theta''}$ has a stratification by 
\[ \sY^{\ch\lambda-\ch\nu, \ch\theta'} \ttimes (\msf S^{\ch\nu}\cap \Gr^{\ch\theta''}) \ttimes C,\]
where $\ch\nu$ ranges over weights of $V^{\ch\theta''}$ (in particular, $\ch\nu\ge \ch\theta''$). 
By Corollary~\ref{cor:Mcover}(ii), if
$\sY^{\ch\lambda-\ch\nu,\ch\theta'}$ maps to a connected component $M$ of 
$\sM^{\ch\theta'}_X$, then $\sY^{\ch\lambda-\ch\theta'', \ch\theta'}$ must map to the
same connected component. 
Therefore $Z^{?,\ch\lambda,\ch\theta',\ch\theta''}$ is irreducible with 
dense open stratum
$\sY^{\ch\lambda-\ch\theta'',\ch\theta'} \times C$. 
Therefore 
$\dim Z^{?,\ch\lambda,\ch\theta',\ch\theta''} = \dim(\sY^{\ch\lambda-\ch\theta'',\ch\theta'})+1
\ge \len(\ch\lambda-\ch\theta'-\ch\theta'') + 1$, as desired.
\end{proof}

This completes the proof of Theorem~\ref{thm:heckeIC}.





\bibliographystyle{amsalpha}
\bibliography{sphericalL}
